\subsection{A bounded subsolution and the limiting equation}
\label{sec:bounded-subsolution-ball-equation}
We can now identify the canonical limit. A bounded local subsolution supplies a zero-boundary barrier near the support of the prescribed measure. Comparison with a strict perturbation of that barrier gives a bound for the limit and recovers its boundary values.

\begin{lemma}\label{lem:capacity-controlled-measure-weights}
Let $\nu_j,\nu$ be finite positive measures supported in a fixed compact set $K_0\Subset B$, with uniformly bounded masses, and suppose $\nu_j\to\nu$ weakly. Assume there is a nondecreasing function $\eta:[0,+\infty]\to[0,+\infty]$ with $\eta(t)\to0$ as $t\downarrow0$ such that, for every Borel $E\subset B$,
\[
 \nu_j(E),\ \nu(E)\leq\eta(\operatorname{Cap}_B^*(E)).
\]
Let $f_j,f$ be real Borel functions with a common finite bound, let $f$ be quasi-continuous near $K_0$, and suppose $f_j\to f$ in outer capacity on $K_0$. Then
\[
 \int_Bf_j\,\mathrm d\nu_j\longrightarrow\int_Bf\,\mathrm d\nu.
\]
For every $c\in\mathbb R$ with $\nu(\{f=c\})=0$, also
\begin{equation}\label{eq:capacity-controlled-moving-level-sets}
 \nu_j(\{f_j<c\})\longrightarrow\nu(\{f<c\}).
\end{equation}
\end{lemma}
\begin{proof}
Write $C$ for the common function bound and $M$ for the mass bound. First keep the weight $f$ fixed. On a compact neighborhood of $K_0$, remove an open $G$ of arbitrarily small capacity so that $f$ is continuous on the remaining compact set. Its continuous extension from \cref{sec:bt-quasicontinuity}, multiplied by a compact cutoff equal to one near $K_0$, gives $g\in C_c(B)$ with $|g|\leq C$ and $g=f$ on $K_0\setminus G$. Weak convergence applies to $g$, and the two replacement errors are at most $2C\nu_j(G)$ and $2C\nu(G)$. Let the capacity of $G$ tend to zero. This proves convergence for the fixed quasi-continuous weight.

For the varying weights and $\delta>0$,
\[
 \int_B|f_j-f|\,\mathrm d\nu_j
 \leq\delta M+2C\eta\left(
       \operatorname{Cap}_B^*(K_0\cap\{|f_j-f|>\delta\})\right).
\]
Let $j\to\infty$, then $\delta\downarrow0$, and combine with the fixed-weight result.

For the level-set assertion choose continuous piecewise linear functions $\sigma^-_s,\sigma^+_s$ with values in $[0,1]$ such that
\[
 \mathbf1_{(-\infty,c-s)}\leq\sigma^-_s
 \leq\mathbf1_{(-\infty,c)}\leq\sigma^+_s
 \leq\mathbf1_{(-\infty,c+s]},\qquad s>0.
\]
For example the lower ramp falls linearly from one at $c-s$ to zero at $c$, and the upper one from one at $c$ to zero at $c+s$. Their compositions with $f_j$ converge in capacity to the corresponding compositions with $f$, whose quasi-continuity is preserved. Apply the first assertion to these compositions and squeeze. As $s\downarrow0$, dominated convergence and the null level set give \eqref{eq:capacity-controlled-moving-level-sets}.
\end{proof}

\begin{corollary}\label{cor:energy-controlled-fixed-measure-limit}
Suppose $\varphi_j\in\mathcal E_0(B)$, $\sup_je_2(\varphi_j)<\infty$, and $\varphi_j\to u$ in local volume $L^1$, with the canonical nonpositive psh limit. If $\nu$ is a finite positive Borel measure and $\nu\leq(\ddc w)^n$ for a bounded $w\in\PSH(B)$, then $u\in L^2(\nu)$ and $\varphi_j\to u$ in $L^1(\nu)$.
\end{corollary}
\begin{proof}
The moment theorem with $p=2$ gives uniform $L^2(\nu)$ bounds for $\varphi_j$. For fixed $T$, the bounded truncations converge in $L^1(\nu)$ by \cref{lem:bounded-psh-fixed-measure-continuity}, and their squares converge in integral. Letting $T\to\infty$ gives the same $L^2$ bound for $u$. The error between a negative function and its truncation has $L^1(\nu)$ norm at most $T^{-1}$ times its squared $L^2(\nu)$ norm. Apply this to $\varphi_j,u$, first taking $j\to\infty$ at fixed $T$ and then $T\to\infty$.
\end{proof}

\begin{theorem}\label{thm:compact-bounded-subsolution-ball-dirichlet}
Let $\mu$ be a finite positive measure of compact support in a ball $B\subset\mathbb C^n$, $n\geq1$. If $\mu\leq(\ddc w)^n$ for a bounded $w\in\PSH(B)$, there is a unique bounded psh function $U$ with
\[
 (\ddc U)^n=\mu,\qquad
 \lim_{B\ni z\to\zeta}U(z)=0\quad(\zeta\in\partial B).
\]
It belongs to $\mathcal E_0(B)$ and is the pointwise envelope of bounded psh subsolutions with upper boundary limits at most zero. The entire canonical sequence in \cref{prop:canonical-grid-energy-control} converges to $U$ in $L^1(B)$ and locally in outer capacity, and satisfies
\[
 \int_B|U_j-U|\,\mathrm d\mu
       +\int_B|U_j-U|\,\mathrm d\mu_j\longrightarrow0.
\]
\end{theorem}
\begin{proof}
If $\mu=0$, the homogeneous zero-boundary solution and all its canonical approximants are zero. The maximum principle identifies the stated envelope. Write $m=\mu(B)>0$ for the remaining case.

\emph{Step 1. A zero-boundary barrier.}
Write $B=B(a,R)$, $\rho(z)=|z-a|^2-R^2$ and $w_0=w-\sup_Bw-1$. Then $-M-1\leq w_0\leq-1$, where $M$ is the oscillation of $w$. Choose $A>0$ so large that $A\rho<-M-2$ on a neighborhood of $\operatorname{supp}\mu$. The function
\[
 W=\max\{w_0,A\rho\}
\]
equals $w_0$ on that neighborhood and equals $A\rho$ on a boundary collar. It is bounded, negative and has zero boundary values; maximum locality gives $(\ddc W)^n\geq\mu$.

Its total mass is finite, since the collar measure is smooth and the remaining interior measure is locally finite. Apply the compact-support polynomial-current identity in the proof of \cref{lem:zero-boundary-truncation-mass} to $W$ and $A\rho$. Their difference has compact support, so testing its $\ddc$ with a cutoff equal to one there gives
\[
 \int_B(\ddc W)^n=A^n\int_B(\ddc\rho)^n.
\]
The same argument gives finite mass for $W_\varepsilon=W+\varepsilon\rho$, $\varepsilon>0$, whose collar is $(A+\varepsilon)\rho$. Positivity of mixed products gives
\begin{equation}\label{eq:local-subsolution-strict-barrier}
 \nu_\varepsilon:=(\ddc W_\varepsilon)^n
                       \geq\mu+\varepsilon^n(\ddc\rho)^n.
\end{equation}

\emph{Step 2. Comparison along a canonical subsequence.}
Choose an $L^1(B)$-convergent subsequence of $U_j$, with limit $u$. The preceding results give local outer-capacity convergence and a uniform second energy bound $E_2$. By the corollary, $U_j\to u$ in $L^1(\nu_\varepsilon)$ for each fixed $\varepsilon>0$.

The grid measures have supports in a fixed $K_0\Subset B$. They have a common capacity modulus by \eqref{eq:energy-measure-capacity-modulus}:
\[
 \mu_j(E)\leq T^n\operatorname{Cap}_B^*(E)+E_2T^{-2}.
\]
The original measure satisfies $\mu(E)\leq M^n\operatorname{Cap}_B^*(E)$. Define $\eta(t)$ as the supremum of their masses on Borel sets of outer capacity at most $t$. First choosing $T$ large and then making the capacity small shows that this common nondecreasing modulus tends to zero. Thus the weight lemma applies to $\mu_j,\mu$.

Fix $\varepsilon>0$. Choose $0<\delta<1$ such that the level set $\{u-W_\varepsilon=-\delta\}$ has zero mass for both $\mu$ and $\nu_\varepsilon$. Only countably many levels are excluded, since both measures are finite. For $T>\|W_\varepsilon\|_\infty+2$, the inequalities
\[
 U_j<W_\varepsilon-\delta,\qquad u<W_\varepsilon-\delta
\]
are unchanged when $U_j,u$ are replaced by their truncations at $-T$. The bounded differences between those truncations and $W_\varepsilon$ are quasi-continuous and converge in capacity. Hence \eqref{eq:capacity-controlled-moving-level-sets} gives
\[
 \mu_j(\{U_j<W_\varepsilon-\delta\})
                       \longrightarrow\mu(\{u<W_\varepsilon-\delta\}).
\]
The fixed-measure $L^1(\nu_\varepsilon)$ convergence gives the corresponding convergence for $\nu_\varepsilon$. To see this directly, outside $\{|U_j-u|>s\}$ the two sublevel indicators can differ only in $\{|u-W_\varepsilon+\delta|\leq s\}$. The first set has mass tending to zero; let $s\downarrow0$ in the second, using the null level set.

The comparison set for $U_j$ and $W_\varepsilon-\delta$ is relatively compact because $U_j$ has zero boundary limit and $W_\varepsilon-\delta\leq-\delta$. Bounded comparison gives
\[
 \nu_\varepsilon(\{U_j<W_\varepsilon-\delta\})
                       \leq\mu_j(\{U_j<W_\varepsilon-\delta\}).
\]
Pass to the limit. By \eqref{eq:local-subsolution-strict-barrier}, the set $\{u<W_\varepsilon-\delta\}$ has zero Euclidean volume. Thus $u\geq W_\varepsilon-\delta$ almost everywhere, and canonical psh representatives preserve this ordering at every point. Let $\delta\downarrow0$ through admissible levels, then $\varepsilon\downarrow0$. We obtain $W\leq u\leq0$.

The barrier makes $u$ bounded and gives its zero boundary limits. For every $T>\|u\|_\infty$, its truncation is $u$ itself. The estimate \eqref{eq:canonical-grid-truncated-measure-limit}, followed by $T\to\infty$, proves $(\ddc u)^n=\mu$. Its mass is $m$, so $u\in\mathcal E_0(B)$.

\emph{Step 3. Uniqueness and the full sequence.}
Let $v$ be a bounded psh subsolution with upper boundary limits at most zero, and put $v_\varepsilon=v+\varepsilon\rho$. Then $(\ddc v_\varepsilon)^n\geq\mu+\varepsilon^n(\ddc\rho)^n$. Comparison with $u$ gives
\[
 \mu(\{u<v_\varepsilon\})
       +\varepsilon^n\int_{\{u<v_\varepsilon\}}(\ddc\rho)^n
          \leq\mu(\{u<v_\varepsilon\}).
\]
The two $\mu$ integrals are finite. Their cancellation makes the strict set volume-null, so canonical representatives give $u\geq v_\varepsilon$ everywhere. Let $\varepsilon\downarrow0$. Since $u$ is itself a subsolution, this proves the pointwise envelope assertion. Interchanging two bounded solutions proves uniqueness; call the solution $U$.

Every subsequence of the canonical sequence has a further $L^1(B)$-convergent subsequence. The argument above identifies every such limit with $U$; \cref{thm:canonical-grid-self-integral-convergence,thm:canonical-grid-capacity-convergence} give its energy and capacity convergences. If any asserted full-sequence convergence failed, a subsequence separated from its stated limit would have such a further convergent subsequence, a contradiction. This proves all the assertions.
\end{proof}
\index{Dirichlet problem!bounded subsolution and compact measure}
